\documentclass[11pt,oneside,openany]{book}
\usepackage[T1]{fontenc}
\usepackage{lmodern}
\usepackage[letterpaper,inner=1.12in,outer=1.12in,top=1.0in,bottom=1.0in]{geometry}
\usepackage{amsmath,amssymb,amsthm,mathtools,booktabs,longtable,microtype,xcolor}
\usepackage{hyperref,xurl,bookmark,etoolbox}
\definecolor{navy}{HTML}{17304A}
\definecolor{teal}{HTML}{1B615F}
\hypersetup{colorlinks=true,linkcolor=navy,urlcolor=teal,citecolor=teal,
  pdftitle={The Relative SU(2) BFSS Quantum System, Edition 0.10.0},
  pdfauthor={Evan Thomas Kotler and the Internal mathematical research collaboration},
  pdfsubject={Closed supercharges and the zero-energy existence problem}}
\urlstyle{same}
\setcounter{tocdepth}{0}
\setcounter{secnumdepth}{1}
\makeatletter
\patchcmd{\l@chapter}{\vskip 1.0em \@plus\p@}{\vskip .42em \@plus\p@}{}{}
\patchcmd{\l@part}{\addvspace{2.25em \@plus\p@}}{\addvspace{1.15em \@plus\p@}}{}{}
\makeatother
\setlength{\emergencystretch}{2em}
\renewcommand{\arraystretch}{1.18}

\newtheorem{theorem}{Theorem}[chapter]
\newtheorem{proposition}[theorem]{Proposition}
\newtheorem{lemma}[theorem]{Lemma}
\theoremstyle{definition}
\newtheorem{definition}[theorem]{Definition}
\newcommand{\E}{\mathcal E}
\newcommand{\F}{\mathcal F}
\newcommand{\Ltwo}{\mathcal H_{\mathrm{amb}}}
\newcommand{\Core}{\mathcal D}
\newcommand{\PhysCore}{\mathcal D_{\mathrm{phys}}}
\newcommand{\ip}[2]{\langle #1,#2\rangle}
\newcommand{\norm}[1]{\lVert #1\rVert}
\newcommand{\op}{\mathrm{op}}
\newcommand{\dom}{\operatorname{Dom}}
\newcommand{\ran}{\operatorname{ran}}
\newcommand{\tr}{\operatorname{tr}}
\newcommand{\Spin}{\operatorname{Spin}}
\newcommand{\SU}{\operatorname{SU}}
\newcommand{\SO}{\operatorname{SO}}
\newcommand{\Gram}{G}
\newcommand{\R}{\mathbb R}
\newcommand{\C}{\mathbb C}
\newcommand{\A}{\overline Q}
\newcommand{\snap}{ec55da1d6e6ffe50fedc1ea722f63eb470fae021}
\newcommand{\repo}{https://github.com/etblink/math/blob/\snap}
\newcommand{\run}[1]{\href{https://github.com/etblink/math/actions/runs/#1}{run \##1}}
\newcommand{\src}[2]{\href{\repo/#1}{#2}}
\newcommand{\evidence}[1]{\par\smallskip\noindent\begingroup\small\color{navy}\textbf{Evidence.}\enspace #1\par\endgroup}
\newcommand{\chaplead}[1]{\noindent\textit{#1}\par\bigskip}

\begin{document}
\frontmatter
\begin{titlepage}
\thispagestyle{empty}
\noindent\textcolor{teal}{\large\bfseries Internal research monograph}\par
\vspace{1.25in}
\noindent{\color{teal}\rule{\textwidth}{1.1pt}}\par
\vspace{0.28in}
{\color{navy}\noindent\Huge\bfseries The Relative $\SU(2)$\\[0.15em] BFSS Quantum System\par}
\vspace{0.32in}
{\noindent\LARGE Closed Supercharges and the\\[0.12em] Zero-Energy Existence Problem\par}
\vfill
{\noindent\Large Edition 0.10.0\par}
\smallskip
{\noindent\large\bfseries Evan Thomas Kotler\par}
{\noindent Internal mathematical research collaboration\par}
{\noindent 27 November 2026, Pacific time\par}
\vspace{0.48in}
\end{titlepage}

\chapter*{Preface}
\addcontentsline{toc}{chapter}{Preface}
This edition studies one spectral question on one specified operator. Does the original relative $\SU(2)$ matrix quantum mechanics have a normalizable zero-energy state? The question takes us from the literal sixteen supercharges to their closed graphs and from a source-defined trial vector to a proposed lower bound on its ground-subspace projection. The argument exposes where the mathematical work stands and precisely what evidence would change the conclusion.

The book assumes graduate linear algebra and basic Hilbert space theory. Chapters~\ref{ch:question}--\ref{ch:charges} establish the notation and the domain chain. Chapter~\ref{ch:square} gives the conventional proof route from the qualified graph pairing to the common square. Chapters~\ref{ch:spin}--\ref{ch:certificate} follow the positive-regulator estimate through the exact trial calculation. Chapters~\ref{ch:local}--\ref{ch:comparison} cover independent spectral, topological and empirical questions. Appendix~\ref{app:evidence} gives the claim index. Appendix~\ref{app:crosswalk} locates material from Edition 0.9.1.

The repository research cutoff is \href{https://github.com/etblink/math/commit/\snap}{\texttt{ec55da1d6e6f}}. The latest qualification in the common-square sequence is CSQ-F0 R3. Each source receipt names its own tested commit. The book uses ordinary mathematical proofs alongside source accounts, finite computational checks, and empirical benchmarks. Their evidence labels apply at the level of the stated claim. The main spectral conclusion remains conditional at this cutoff.

\tableofcontents
\mainmatter
\part{The existence problem}
\chapter{A Zero-Energy State in the Relative Model}\label{ch:question}
\chaplead{The spectral target is a vector in the kernel of the original nonnegative Hamiltonian. An explicit trial vector can establish existence by having positive projection onto that kernel.}

\section{Which problem is being solved?}
The relative $\SU(2)$ system is a matrix quantum mechanics with nine spatial labels and three adjoint color labels. Its bosonic configuration space is $\mathbb R^{3\times9}$ and its fermionic fiber is the complex Fock space of twenty-four modes. Gauge invariance selects a closed Hilbert subspace $\E$ of the ambient square-integrable wavefunctions. The Hamiltonian $H$ used throughout is the normalized adjoint square of the original closed charge column $T$ on that subspace.

Banks, Fischler, Shenker and Susskind proposed matrix quantum mechanics as part of a large-$N$ description of M-theory \cite{bfss}. This book concerns the finite relative model. The physical subspace here is defined by the original gauge symmetry. Identifying this space and proving a spectral statement on it are mathematical tasks. Empirical interpretation requires an additional map to physically specified observables.

For $\psi\in \E$ with $\norm{\psi}=1$, write
\begin{equation}
 F(\psi)=\norm{P_{\ker H}\psi}^{2}.\label{eq:fidelity}
\end{equation}
If $F(\psi)>0$, the vector $P_{\ker H}\psi/\sqrt{F(\psi)}$ is a normalizable zero mode. The spectral projection is supported at the point zero, so its range is in $\dom H$ and $H$ vanishes there. The statement concerns existence. Any assertion about multiplicity, a gap, or higher point spectrum needs further information.

\section{Why small energy leaves a question}
A trial energy estimates a first moment of the spectral measure of $\psi$. It cannot generally lower-bound the atom at zero. For example, on $L^2(0,1)$ take $H$ to be multiplication by $s$, and $\psi_\varepsilon=\varepsilon^{-1/2}\mathbf 1_{(0,\varepsilon)}$. Its energy is $\varepsilon/2$ and $\ker H=\{0\}$. The example explains why a sequence of small expectation values on its own leaves the zero-mode question open.

The positive-regulator route controls the part of $\psi$ orthogonal to $\ker H$ using its weighted charge outputs. With $r=\norm{X}$ and $\delta>0$, the target inequality is
\begin{equation}
 F(\psi)\ge 1-\kappa q_\delta[\psi],
 \qquad
 q_\delta[\psi]=\frac1{16}\sum_{\alpha=1}^{16}
      \norm{\sqrt{r^2+\delta^2}\,\A_\alpha\psi}^{2}.
 \label{eq:target}
\end{equation}
The source Hardy estimate in Chapter~\ref{ch:fidelity} gives the proposed $\kappa=(36+8\sqrt{14})/25$. The bound uses the physical Spin(9) action on charge outputs, the common square, and a weighted-domain trial singlet. The finite computation must measure the original charge form on that trial. Chapter~\ref{ch:certificate} states the complete implication with those premises.

\section{The evidence at the cutoff}
The original normalized closed Hamiltonian is available, along with the equality of the individual and joint form domains and essential self-adjointness of the literal invariant-core charges. CSQ-F0 R3 qualifies equality of complex component pairings on the original closed joint graph. The transport of these pairings to the native normalized form and the identification of the natural squares have an analytical route. Physical Spin(9) covariance, the Hardy implication and the identification of the finite trial data have their own source obligations.

The 104-function construction supplies the worked finite example. Its archived rational endpoint at $\delta=1$ is below $3/8$ under the source identification. For $\kappa<66/25$, the product of these two rational upper bounds is $99/100$. This explains why even the conservative threshold would suffice. Smaller positive regulators give wider conditional margins, and reduced witnesses offer smaller independent replay tasks. The compactness and fixed-ball route in Chapter~\ref{ch:local} is independent of this numerical margin.

\section{How statements are used}
A source qualification concerns exact declarations, hypotheses and tested bytes. A conventional proof in this edition exposes the deductions a reader can check on paper. A finite computational audit tests a defined finite set of operations, usually with exact arithmetic. An empirical benchmark records an observable that a proposed physical map would have to explain. The evidence index gives separate entries for these claims, with links to the original artifacts where available.

The governing research question is what follows from the evidence. The source mathematics may produce a constructive result, an obstruction, or a sharper question. Chapter~\ref{ch:comparison} applies the same standard to the Hypothesis of Everything and proposed links to particle physics.
\evidence{\src{research/methodology/TRUTH_FIRST_FOUNDATIONAL_RESEARCH_CHARTER_2026_11_18.md}{Truth-first research charter}.
The source and analytical statuses are indexed in Appendix~\ref{app:evidence}.}

\chapter{The Original Carriers and Sixteen Charges}\label{ch:source}
\chaplead{The normalization of the color basis, Majoranas, fermion field and charge column fixes the factors used in every subsequent estimate.}

\section{Bosonic coordinates and color}
Use $X_{ai}=x_i^a$ for three color rows and nine spatial columns. Set
\[
 r^2=\tr(XX^{\mathsf T}),\qquad
 \Gram=XX^{\mathsf T},\qquad
 s=\tr\Gram,\quad u=\tfrac12((\tr\Gram)^2-\tr(\Gram^2)),\quad v=\det\Gram .
\]
The letter $u$ is an invariant in formulas involving $(s,u,v)$ and is a Hilbert vector in operator-domain formulas. The source uses zero-based indices for the finite types. We use one-based spatial and spin labels in displayed formulas.

Let $e_a=\sigma_a/\sqrt{2}$ be the normalized Hermitian Pauli color basis. Its trace pairing is orthonormal. For the source bracket $[E,F]_{\rm s}=-i(EF-FE)$, the structure coefficients are $f_{abc}=\sqrt{2}\,\epsilon_{abc}$ in the oriented Pauli order. Indeed $[\sigma_b,\sigma_c]=2i\sum_a\epsilon_{abc}\sigma_a$, and the two input factors $1/\sqrt{2}$ give $[e_b,e_c]_{\rm s}=\sum_a\epsilon_{abc}\sigma_a$. This convention fixes the crossover to the ordinary cross product.

\section{The fermionic fiber}
The real symmetric $16\times16$ matrices $\gamma_i$ obey
\[
 \gamma_i\gamma_j+\gamma_j\gamma_i=2\delta_{ij}I,\qquad
 \gamma^{ij}=\tfrac12(\gamma_i\gamma_j-\gamma_j\gamma_i).
\]
They act on a spatial spin label. The forty-eight original Hermitian fermions $\theta_{\alpha a}$ act on $\F$, with anticommutator
\[
 \{\theta_{\alpha a},\theta_{\beta b}\}=\delta_{\alpha\beta}\delta_{ab}I.
\]
It is convenient to use $t_{\alpha a}=\sqrt{2}\,\theta_{\alpha a}$, which square to $I$. In a paired mode with creation and annihilation $c_j^\dagger,c_j$ satisfying the complex CAR, the two normalized Majoranas are $c_j+c_j^\dagger$ and $i(c_j^\dagger-c_j)$ up to the oriented sign convention of the chosen pairs. The computational packing uses the actual source pair order and signs, which are part of its isometry receipt.

The Fock space is $\bigwedge^\bullet\C^{24}$ and $\dim_\C\F=2^{24}$. The multiplication algebra generated by all normalized Majoranas is $\operatorname{End}_\C\F$. One can see surjectivity directly: the number operators give the vacuum projection $\prod_j(1-c_j^\dagger c_j)$, and creation and annihilation words on its two sides give matrix units between occupation basis vectors. The dimension of the complex Clifford algebra for forty-eight generators is $2^{48}$, equal to the dimension of this endomorphism algebra. The full algebra has scalar commutant. Later uses on actual gauge states and accessible observables require their own constraints.

\evidence{\src{lean/OAI/MathematicalPhysics/BFSS/Core.lean}{Original source core}.
The theorem atlas records the carrier and original CAR results. The packed transport used for the trial is indexed in Chapter~\ref{ch:witness}.}

\section{Smooth core and gauge action}
Let $\Ltwo=L^2(\R^{27},dX;\F)$ with inner product conjugate-linear in its first argument. The compact smooth core $\Core=C_c^\infty(\R^{27};\F)$ embeds injectively in $\Ltwo$: a continuous function that vanishes almost everywhere with respect to full-support Lebesgue measure vanishes everywhere.

Write $R_g$ for the adjoint color matrix and $V_g$ for the actual fermionic unitary of the source gauge action. In these conventions the gauge action on sections is
\[
 (U_g f)(X)=V_g f(R_g^{-1}X).
\]
The adjoint color matrices lie in $\SO(3)$ and the adjoint map from $\SU(2)$ is onto. Orthogonality makes the bosonic change of variables measure preserving. The physical Hilbert space and core are
\[
 \E=\bigcap_g\ker(U_g-I),\qquad
 \PhysCore=\{f\in\Core:f(R_gX)=V_g f(X)\ \forall g,X\}.
\]
Boundedness of each $U_g$ makes the fixed space closed. For compact gauge group let $P_G f=\int U_gf\,dg$ with Haar probability measure. Averaging is contractive and fixes physical vectors. On compact smooth inputs the union of the rotated supports is bounded, and derivatives can be taken under the integral. Hence $P_G\Core\subset\PhysCore$. Applying $P_G$ to ambient smooth approximants proves density of the physical core in $\E$. The corresponding source Haar operation is used again at graph level.

\section{Nontriviality and the Gauss test}
An original nonzero gauge-invariant Fock vector multiplied by a nonzero smooth radial compact bump is a physical normalizable vector. The radial factor is invariant under the adjoint orthogonal color action. The norm is the product of its scalar $L^2$ norm and the fixed fiber norm, so positivity is immediate. The construction establishes that $\E$ is nonzero. Its differential energy is a separate calculation.

The source also has a historical maximal-domain Gauss characterization. On the ambient Hilbert space, the joint kernel of the actual Gauss generators at their maximal domains equals the fixed space of the genuine $\SU(2)$ action. That identification is stronger than the core assertion obtained by differentiating gauge invariance of smooth functions. Its historical receipt and exact domain interface are retained in the atlas.
\evidence{The historical radial construction and maximal Gauss equivalence are atlas G-02 and G-04. The original invariant fiber vacuum is G-03.}

\section{The literal differential operators}
For each $\alpha$ the source supercharge is
\begin{align}
 Q_\alpha f={}&-i\sum_{i,a,\beta}(\gamma_i)_{\alpha\beta}
           \theta_{\beta a}\,\partial_{ia} f \nonumber\\
 &+\frac12\sum_{i,j,a,b,c,\beta}f_{abc}(\gamma^{ij})_{\alpha\beta}
           \theta_{\beta a}x_i^b x_j^c f .
 \label{eq:sourcecharge}
\end{align}
The spatial pair in the second line is ordered. Equivalently, writing $X_i$ for the $i$th color column and using the normalized Majoranas gives the unordered-pair expression
\[
 Q_\alpha=-\frac{i}{\sqrt2}\sum_{i,a,\beta}
       (\gamma_i)_{\alpha\beta}t_{\beta a}\,\partial_{ia}
   +\sum_{i<j,a,\beta}(\gamma_i\gamma_j)_{\alpha\beta}
       (X_i\times X_j)_a\,t_{\beta a}.
\]
The $i<j$ pair occurs twice in the ordered sum, $f_{abc}$ supplies the factor $\sqrt2$, and $\theta=t/\sqrt2$. These factors explain the displayed coefficient of one for the polynomial term.

Each $Q_\alpha$ takes smooth compact inputs to smooth compact outputs. Its first-order coefficients have the form $-i$ times constant Hermitian matrices, and its zeroth-order field is pointwise Hermitian. Integration by parts gives $\ip{Q_\alpha f}{g}=\ip{f}{Q_\alpha g}$ on the ambient core. Gauge covariance makes the invariant core stable under the charges. For the derivative term this uses the orthogonal chain rule and the gauge action on the Majoranas. For the polynomial term it uses preservation of the color bracket under conjugation.

The source quadratic form is
\[
 q_{\rm core}[f]=\frac1{16}\sum_{\alpha=1}^{16}\norm{Q_\alpha f}^2 .
\]
This charge normalization carries through to the adjoint-square Hamiltonian. Formal symmetry and the finite sums are the starting data for the closed-operator work.

\part{From supercharges to the Hamiltonian}
\chapter{The Native Hamiltonian and the Massless Interaction}\label{ch:ham}
\chaplead{The original closed charge graph determines the operator, the form domain and the test identity required to recognize its action.}

\section{The original joint graph}
Embed the physical test graph
\[
 G_0=\{(f,(Q_\alpha f)_\alpha):f\in\PhysCore\}
       \subset\Ltwo\times \Ltwo^{16}
\]
in the ambient product. Its graph closure is a graph: if $f_n\to0$ and $Q_\alpha f_n\to y_\alpha$ for every $\alpha$, then for any ambient core vector $g$,
\[
 \ip{y_\alpha}{g}=\lim_n\ip{Q_\alpha f_n}{g}
                =\lim_n\ip{f_n}{Q_\alpha g}=0.
\]
Density of the ambient core gives $y_\alpha=0$. The finite product topology and the finite Hilbert-sum norm give the same closure, and the latter specifies the inner product for the adjoint. Identifying the input with its physical subtype and the output with its Hilbert-sum coordinates gives the original closed operator
\[
 T:D(T)\subset\E\longrightarrow\Ltwo^{16}.
\]
Since $\PhysCore$ is dense in $\E$, $T$ is densely defined. At source level the graph is pulled back by the actual physical inclusion and the canonical Hilbert-sum coordinate map. This construction keeps the original charge graph as the reference object.

The all-space graph, defined via $\Core$ in place of $\PhysCore$, admits Haar averaging on both entries. Compatibility of the smooth Haar action with the full charge column extends to the closed graph by boundedness and approximation. For an all-space graph point with physical input, its Haar average has that same input and therefore the same output, by single-valuedness. The averaged approximating core vectors lie in $\PhysCore$. This explains the physical-slice equality of the two graphs that the original source gate establishes.
\evidence{AY gives the native physical Hilbert-sum graph and the exact coordinate transport. UGC4 R3 qualifies the restricted/full physical-input graph equality. \run{45869599441}.}

\section{Domain, action and weak testing}
For a closed densely defined $T$, its positive self-adjoint adjoint square is defined on
\[
 D(T^\ast T)=\{f\in D(T):Tf\in D(T^\ast)\}.
\]
For the normalized Hamiltonian set
\[
 H=\tfrac1{16}T^\ast T,\qquad
 h(f,g)=\tfrac1{16}\ip{Tf}{Tg},\qquad D(h)=D(T).
\]
The form norm $\sqrt{\norm{f}^2+h[f]}$ is equivalent to the graph norm of $T$. The graph is closed, so the form is closed. The scalar $1/16$ leaves the adjoint-square domain unchanged and determines the action.

For $f\in D(T)$ and $w\in\E$, the Hilbert adjoint definition gives
\begin{equation}
 \bigl(\forall g\in D(T),\ \ip{w}{g}=h(f,g)\bigr)
 \quad\Longleftrightarrow\quad f\in D(H),\ Hf=w.
 \label{eq:weakH}
\end{equation}
To see the normalization, the test equation is equivalent to
$\ip{Tf}{Tg}=\ip{16w}{g}$ for all $g$. This is precisely $Tf\in D(T^\ast)$ and $T^\ast Tf=16w$. The weak identity also shows why replacing the core by the closed form domain matters.

A vector in $D(T)$ with $Tf=0$ satisfies the weak equation with $w=0$. Conversely, $Hf=0$ implies $h[f]=0$. Thus $\ker H=\ker T$. This proves a kernel equivalence on the physical carrier and leaves its dimension as a spectral question.
\evidence{The exact normalized weak representative, self-adjointness, and closedness are in the BE source interface. Appendix~\ref{app:evidence} gives the declaration map.}

\section{Historical comparison of Hamiltonians}
The original research archive also defines a historical physical operator via a matched carrier and form. Equality of operators under that identification needs equality of domains and outputs. The H-03 historical result provides a comparison at its tested source snapshot and explicit matching hypotheses. Changes to the port's imports and local instance expression are tracked in the historical atlas. For the rest of this book $H$ is the native normalized operator of \eqref{eq:weakH}; it has the independent native source receipt indexed by E08-BF.
\evidence{H-03 in the \src{contributions/bfss-kotler/internal-monograph/THEOREM_ATLAS.json}{historical theorem atlas}; E08-BF \run{43498649514}.}

\section{The massless field and energy}
Let $Y(X)$ be the unnormalized Hermitian fermionic field of the original massless source, linear in $X$. As a wavefunction multiplier the natural maximal domain consists of $f\in\Ltwo$ such that $Yf\in\Ltwo$. To see closedness, suppose $f_n\to f$ and $Yf_n\to g$ in $L^2$. A subsequence converges almost everywhere in both coordinates. Pointwise multiplication by the finite matrix $Y(X)$ gives $g=Yf$ almost everywhere. On the physical subspace the source covariance transports both membership in this domain and the output action. In particular smooth compact vectors lie in the multiplier domain.

In the concrete Pauli normalization the source-matched ordinary energy on the invariant smooth core is
\begin{equation}
 q_{\rm core}[f]=\tfrac12\norm{\nabla f}^2+
       \int \left(u(X)\norm{f(X)}^2+
       \Re\ip{f(X)}{Y(X)f(X)}\right)\,dX .
 \label{eq:energy}
\end{equation}
Here $u(X)=\sum_{a<b}\norm{X^a\wedge X^b}^2$ is the second elementary symmetric polynomial of $\Gram$. The original massless charge identity specializes the deformed source at its undeformed parameter values, and the concrete Pauli/Fock pairing gives the stated coefficients. Source results on the multiplier and this energy formula do not describe the maximal distributional domain of every first-order charge.

The linear field satisfies a finite-dimensional bound $\norm{Y(X)}_{\op}\le C r$ for some fixed $C$. The quartic potential vanishes on rank at most one. The lower estimate from \eqref{eq:energy} therefore contains a negative radial term and gives no global positive first-moment bound by itself. Chapter~\ref{ch:local} states the further coercivity estimate that would be sufficient.

The bounded normalized Gram/Fock construction uses a different field. Its normalization gives uniform auxiliary control under the stated Gram hypotheses. Substituting that bound for the original $Y(X)$ estimate would change the energy. The source atlas keeps these two objects distinct.
\evidence{Original energy and the massless multiplier are indexed in A-01, the historical S4B36 receipt, and GM0--GM5. See the \src{contributions/bfss-kotler/internal-monograph/THEOREM_ATLAS.json}{historical atlas}.}

\chapter{Individual Charge Closures on the Physical Space}\label{ch:charges}
\chaplead{A closed realization, its weak graph, and essential self-adjointness of the original core restriction are separate steps.}

\section{The individual and joint domains}
Fix a spin label $a$. The original core equality of squared energies implies
\[
 \norm{Q_a f}^{2}=\frac1{16}\sum_{\alpha=1}^{16}\norm{Q_\alpha f}^2
 \qquad (f\in\PhysCore).
\]
Apply it to $f_n-f_m$ for a sequence in $\PhysCore$. Cauchy convergence of one charge output is equivalent to Cauchy convergence of the entire column. A sequence giving a point in the joint graph therefore gives an individual graph point by taking its $a$th coordinate. An individual graph approximation lifts to a joint graph approximation by completeness of the other fifteen output spaces. The lifted point is unique by single-valuedness of the joint graph.

This argument identifies the input domains rather than assuming that closure commutes with projection. In general, the image of a closed graph under a coordinate projection need not be closed, and individual graph approximations might carry no bound on the remaining coordinates. Here the norm identity gives that bound on differences and makes the two approximations equivalent.

Established gauge preservation at graph level transports each projected output to the physical subspace. The native operators $A_a$ on $\E$ are densely defined, and their domains satisfy
\begin{equation}
 D(A_a)=D(T),\qquad
 \iota A_a f=(Tf)_a\quad(f\in D(T)),                 \label{eq:component}
\end{equation}
where $\iota:\E\hookrightarrow\Ltwo$ is the physical inclusion. The source has a qualified common-domain theorem and exact graph transport interfaces. The displayed output formula describes the analytical combination of those interfaces used for the normalized pairing in Chapter~\ref{ch:square}.
\evidence{ICD-A0 supplies the original ambient graph adapter. ICD-B0 proves the component lift under the equal-norm premise. ICD-C0 specializes the common-domain identification. NAG0 and NAG1 construct the native individual closed operators and qualify the dense common domain. \run{45912386053} and \run{46051319682}.}

\section{Weak graphs and the role of symmetry}
The graph of the adjoint of a densely defined symmetric core restriction has a weak description. On the ambient space and with the inner product convention of Chapter~\ref{ch:source}, $(f,y)$ satisfies the adjoint graph test if
\[
 \ip{y}{g}=\ip{f}{Q_a g}\qquad(g\in\Core).
\]
This is the distributional equation $Q_a f=y$ in the test-pairing sense. Its output is unique by density. At physical level one tests $\PhysCore$, and identifying that weaker test family with the physical graph requires actual Haar transport. Orthogonality of the physical projection and commutation with the source charge are used to extend tests from invariant core vectors to arbitrary core vectors.

One route to the ambient maximal graph is instructive. Let $Q=D+V$, where $D$ has constant first-order coefficients and $V$ is the smooth Hermitian matrix polynomial of the source. For $f,y\in L^2$ with $Qf=y$ distributionally, choose real cutoffs $\chi_R(X)=\chi(X/R)$ that equal one near the origin. Then
\[
 Q(\chi_R f)=\chi_Ry+[D,\chi_R]f,\qquad
 \norm{[D,\chi_R]f}\le C R^{-1}\norm{f}.
\]
Thus both input and output converge in $L^2$. For fixed $R$, $D(\chi_R f)$ is in $L^2$, because the weak output is and $V$ is bounded on the compact support. Convolution with a compactly supported smooth approximate identity commutes with $D$ on these weak pairs. The convolutions converge in $L^2$ together with their $D$ output, and boundedness of $V$ on a single enlarged compact set gives convergence of the $V$ term. A diagonal approximation puts the original weak pair in the core graph closure. The argument avoids requiring all first derivatives separately in $L^2$.

The ambient closure is therefore the maximal weak realization. The gauge projection preserves the ambient graph, which lets the physical weak test characterize the restricted closure. This analytic exposition shows the use of cutoff and mollifier estimates. The native source proof is indexed by its exact sequence of cutoffs, weak-graph lemmas, and physical Haar transfers.
\evidence{\src{research/computational/python-first-20261124/bfss-physical-implication/POST_ICD_ESA_MAXIMAL_GRAPH_MOLLIFIER_AND_SQUARE_AUDIT_2026_11_26.md}{November 26 analytical audit}. LOC6C qualifies the physical weak equivalence. \run{46041935392}.}

\section{Self-adjointness and the literal core}
NAG2 supplies native self-adjointness by comparing the native adjoint graph with the original weak-pairing graph. The comparison determines both domains and outputs. NAG3 identifies the graph closure of the literal original invariant-core operator with that native realization. NAG4 then states the explicit density, symmetry, closability and self-adjoint closure package. For each label $a$ we may use $A_a=\overline{Q_a|_{\PhysCore}}$ as a self-adjoint operator on $\E$.

Essential self-adjointness answers which closure of the core action is canonical. A self-adjoint operator has a self-adjoint square, but identifying it with the specific original $H$ requires agreement of their weak forms on the common domain. The result about the core square is the beginning of that comparison.
\evidence{NAG2 R3 \run{46059907504}, NAG3 R2 \run{46067320278}, NAG4 R2 \run{46074553299}. The artifacts and exact declarations are in the physical-existence atlas.}

\chapter{The Common Square}\label{ch:square}
\chaplead{The qualified theorem is a complex pairing identity on the closed joint graph. The following analytical assembly explains how that identity determines the natural square.}

\section{The complex pairing on core vectors}
For invariant core vectors $f,g$ the source square equality gives
$Q_a(Q_a g)=Q_b(Q_b g)$ in the original core representation. Use source symmetry twice:
\begin{equation}
 \ip{Q_a f}{Q_a g}=\ip{f}{Q_a(Q_a g)}
                 =\ip{f}{Q_b(Q_b g)}
                 =\ip{Q_b f}{Q_b g}.
 \label{eq:corepair}
\end{equation}
This establishes the complex identity directly for two independent inputs. Each pairing uses the actual core embedding. The symmetry theorem used in the first and last equalities tests the ambient smooth core, so the intermediate output is an admissible test.

The source proof of \eqref{eq:corepair} uses the direct symmetry and square-equality chain. The earlier attempt to specialize a more generic pairing lemma in the concrete core proof hit an elaboration recursion bound. The abstract topological transport remains useful and is part of the qualified statement.

\section{Closure in both arguments}
Let $S$ be the range of the core graph map in the ambient product, and let $p_\alpha$ be its continuous projection to the $\alpha$th output coordinate. Fix $x\in S$. The set of $y$ satisfying
\[
 \ip{p_a(x)}{p_a(y)}=\ip{p_b(x)}{p_b(y)}
\]
is closed and contains $S$, hence contains $\overline S$. Now fix $y\in\overline S$. The same reasoning in the first variable gives the equality for every $x\in\overline S$. The graph identity from ICD-A0 identifies $\overline S$ with the original minimal closed joint graph. No limits of unbounded compositions are taken.

\begin{theorem}[Source-qualified closed graph pairing]
For arbitrary points $(u,y)$ and $(v,z)$ in the original ambient closed joint graph of the paired $\SU(2)$ data, and any spin labels $a,b$,
\[
 \ip{y_a}{z_a}=\ip{y_b}{z_b}.
\]
\end{theorem}
\evidence{CSQ-F0 R3, exact candidate \texttt{29289558cb05}, \run{46082764637}.
The independent consumer instantiates both the invariant-core and two-point closed-graph statements.
\src{research/computational/python-first-20261124/bfss-physical-implication/CSQF0_R3_GREEN_QUALIFIED_2026_11_27.md}{Qualification receipt}.}

\section{From a coordinate to the normalized form}
Let $D=D(T)=D(A_a)$, and take $u,v\in D$. The AY graph membership equivalence transports the two native points $(u,Tu)$ and $(v,Tv)$ to ambient graph points. NAG0 and ICD-B0 place each native individual output in the selected coordinate of a joint graph point. Since the joint graph is single-valued for a given input, this yields \eqref{eq:component}. The physical inclusion is isometric. By the closed graph pairing,
\[
 \ip{A_a u}{A_a v}=\ip{(Tu)_a}{(Tv)_a}
 =\ip{(Tu)_b}{(Tv)_b}\qquad\text{for every }b.
\]
The actual inner product on the target of $T$ is the Hilbert sum. In the pinned source this is the identity for $\operatorname{PiLp}$ at exponent two. Summing the last expression over the sixteen spin labels gives
\begin{equation}
 h(u,v)=\tfrac1{16}\ip{Tu}{Tv}=\ip{A_a u}{A_a v} \qquad(u,v\in D).
 \label{eq:normalized-pair}
\end{equation}
The topology used for graph closure and the Hilbert-sum normalization play different roles. The topology makes the coordinate pairings continuous. The Hilbert structure supplies the factor of sixteen.

\section{Why the weak test gives natural-domain equality}
For $u\in D$ and $w\in\E$, \eqref{eq:normalized-pair} identifies the two tests
\[
 \forall v\in D,\quad \ip{w}{v}=h(u,v)
 \qquad\text{and}\qquad
 \forall v\in D,\quad \ip{w}{v}=\ip{A_a u}{A_a v}.
\]
By \eqref{eq:weakH} the first is equivalent to $(u,w)\in\operatorname{graph} H$. By the adjoint definition with the dense common domain the second is equivalent to $(u,w)\in\operatorname{graph}(A_a^\ast A_a)$. Each graph has its first coordinate in $D$, so this proves equality of the entire graphs once the source interfaces are assembled. NAG2 gives $A_a^\ast=A_a$, hence the natural composition $A_a^2$ has exactly the same domain and action as $H$.

\begin{proposition}[Analytical common-square implication]\label{prop:common}
Under the native coordinate and form identifications above, every $A_a^2$ equals the original $H$ as a self-adjoint operator on $\E$. In particular,
\[
 \dom H=\{u\in D:A_a u\in D\},\qquad
 Hu=A_a(A_a u)\quad (u\in\dom H).
\]
\end{proposition}
At the cutoff the source-qualified component pairing is available. The normalized physical pairing assembly and the natural-operator identification have the conventional argument above and await their own strict source qualification. The resolvent proof in Chapter~\ref{ch:fidelity} explicitly uses Proposition~\ref{prop:common}.
\evidence{\src{research/computational/python-first-20261124/bfss-physical-implication/NAG4_INTERVAL_DIRECT_CLOSED_PAIRING_AND_WITNESS_COVERAGE_2026_11_27.md}{Direct pairing and weak-form route}.
The exact AY, NAG0--NAG2 and BE interfaces are listed in Appendix~\ref{app:evidence}.}

\part{From a trial state to a certificate}
\chapter{Spin(9) on Physical States}\label{ch:spin}
\chaplead{The charge outputs of a physical singlet must transform as an actual spinor in the wavefunction Hilbert space.}

\section{Three actions with different carriers}
The real spinor label representation $S$ acts on the sixteen gamma labels. The Fock implementation $V(k)$ acts on fermionic vectors. Its conjugation action $B\mapsto V(k)BV(k)^{-1}$ acts on fermionic operators. Conjugation preserves the adjoint and trace when $V(k)$ is unitary. An automorphism of the operator algebra, even with these properties, still needs an implementing phase convention to determine $V(k)$.

The distinction affects particle decompositions and the Hardy estimate. The endomorphism space carries a tensor-product action of the Fock representation and its dual. The charge outputs in the Hardy estimate lie in the physical vector space. Its spinor intertwiners cannot be supplied by an isotypic calculation in the endomorphism representation.

\section{The source ninth-axis stabilizer}
The source map to $\SO(9)$ restricts to a two-sheeted covering of the orthogonal stabilizer of the last axis. The target stabilizer is continuously isomorphic to $\SO(8)$ by restriction to the perpendicular eight-dimensional space. The accepted source factorization by perpendicular unit-vector Clifford words gives compactness and connectedness of the covering subgroup. This fixes the geometric group and its map. A branching argument additionally needs the action of the classified subgroup on the original physical Fock vectors, with the vector eight-module and central signs identified. This is the representation requirement used by the transverse oscillator route in Chapter~\ref{ch:local}.
\evidence{Source covering \run{42452043377} and target stabilizer \run{42498367504}.}

\section{The proposed phase-fixed vector action}
For a real skew endomorphism $B$ of the Majorana labels, define the quadratic infinitesimal action
\[
 K(B)=\frac14\sum_{p,q}B_{pq}t_p t_q .
\]
The CAR give
\[
 [K(B),t_r]=\sum_p B_{pr}t_p,\qquad
 [K(B),K(C)]=K([B,C]).
\]
Indeed $[t_pt_q,t_r]=2(\delta_{qr}t_p-\delta_{pr}t_q)$ gives the first formula. Applying this derivation to both factors in $K(C)$ and using skewness gives the second as an equality of quadratic expressions. This calculation fixes the Lie algebra action without a scalar ambiguity. Since $K(B)^\ast=-K(B)$, integration of the spatial Lie algebra representation to simply connected $\Spin(9)$ gives a unitary implementation with the prescribed derivative.

For a coordinate plane $(i,j)$ the spin-label generator is $\gamma_i\gamma_j/2$. Its Fock generator is
\[
 K_{ij}=\frac18\sum_{a,\alpha,\beta}
            (\gamma_i\gamma_j)_{\alpha\beta}t_{\alpha a}t_{\beta a}.
\]
For $R(k)\in\SO(9)$ the corresponding action on wavefunctions is the finite Fock unitary times inverse pullback under $X\mapsto XR(k)^{\mathsf T}$. The spatial generator is the sum of the orbital derivative and $K_{ij}$ in the source orientation. The color and spatial rotations act on different label factors, so their infinitesimal actions commute. Exact source gauge factorization and the vector implementation are needed for the resulting global commutation on wavefunctions.

\section{Covariance on the closed form domain}
On the smooth core, the Clifford covariance of $\gamma_i$ and of $\gamma^{ij}$ combines with the orbital chain rule in \eqref{eq:sourcecharge}. The spatial action rotates the charge label by the spinor representation. Write $U(k)$ for the action on $\E$, and use compatible conventions so that for an invariant $\psi$ the map
\begin{equation}
 T_\chi:\C^{16}\longrightarrow \E,\qquad
            T_\chi e_\alpha=A_\alpha\psi                         \label{eq:chi}
\end{equation}
intertwines $S$ and $U$. This statement requires the charge covariance on the actual graph domain and actual physical output space. On core vectors the column rotates by a unitary on the Hilbert sum. Preservation of the core and extension by graph convergence then show invariance of $D(T)$ and of $h$. The weak form \eqref{eq:weakH} gives invariance of $\dom H$ and commutation with $H$ and its bounded shifted resolvents.

This is the analytical route for the original paired source. The monograph uses it as a premise for the fidelity proof. Qualifications for an operator action, an infinitesimal label commutator or a topological stabilizer cover each have their own scope.

\section{Averaging intertwiners instead of classifying all sectors}
Given an intertwiner $S\to D(h)$, approximate the images of a spinor basis by physical core vectors in the form graph norm. Average the resulting finite-column map over $\Spin(9)$ by
\[
 M\longmapsto \int_{\Spin(9)} U(k) M S(k)^{-1}\,dk .
\]
Graph norm invariance makes the averaging contractive, and it fixes an intertwiner. Smoothness and a common compact support bound are preserved by the compact spatial action. This provides smooth intertwiner approximants for the Hardy estimate. It leaves the full isotypic decomposition of $\E$ as a separate question.
\evidence{\src{research/computational/python-first-20261124/bfss-physical-implication/GAUGE_SPIN_GLOBAL_REVIEW_2026_11_25.md}{Original gauge and Spin(9) audit}.
The atlas records the genuine vector and domain covariance as an analytical item.}

\chapter{The Positive-Regulator Hardy Estimate and Fidelity}\label{ch:fidelity}
\chaplead{A Hardy estimate for physical spinor intertwiners gives a lower bound on the kernel projection of a physical trial singlet.}

\section{The weighted graph and the shared kernel}
For $\delta>0$ put $w_\delta(X)=\sqrt{r^2+\delta^2}$ and $B_\delta=w_\delta^{-1}$. Its reciprocal is bounded and injective. Weighted outputs can be defined by the inverse image of $\operatorname{graph} T$:
\[
 (u,z)\in G_\delta\quad\Longleftrightarrow\quad (u,B_\delta z)\in
 \operatorname{graph} T .
\]
The map $(u,z)\mapsto(u,B_\delta z)$ is continuous, so the inverse image is closed. Injectivity makes the weighted output unique. The physical smooth core lies in this domain. At $z=0$ the defining equivalence says that $\ker T_\delta=\ker T$, where $T_\delta$ is the weighted column. Their normalized adjoint squares have this same kernel. The existence argument will use the spectral projection of $H$ throughout.

For $u$ in the weighted domain,
\[
 q_\delta[u]=\tfrac1{16}\norm{T_\delta u}^2
 =q_{\rm aux}[u]+\delta^2 h[u],\qquad
 q_{\rm aux}[u]=\tfrac1{16}\norm{rTu}^2 .
\]
The right side defines a scalar quantity on the trial domain. The argument below only uses $B_\delta$ at positive $\delta$.

\section{Hardy estimate on physical spinor columns}
The source calculation has two parts. First, sum the charge anticommutator for four disjoint spatial planes to obtain the total generator $J$ and the corresponding linear fields $\tau_\alpha$. Put $f_\delta=B_\delta^2$. The source product-rule calculation gives
\begin{equation}
 \sum_{\alpha=1}^{16}\{Q_\alpha,f_\delta\tau_\alpha\}
 = f_\delta J+2i f_\delta^2\sum_{\alpha=1}^{16}\Xi_\alpha\tau_\alpha .
 \label{eq:anticomm}
\end{equation}
The $\Xi_\alpha$ come from the radial multiplier differentiation. The finite contraction identities in the audit are
\[
 \sum_\alpha\norm{\Xi_\alpha(X)v}^2=8r^2\norm{v}^2,\qquad
 \sum_\alpha\norm{\tau_\alpha(X)v}^2
   =\frac{61r^2-45\sum_a X_{a9}^2}{1568}\norm{v}^2 .
\]
Both coefficients refer to the original charge normalization. They enter the physical argument after compression to an intertwiner.

Second, on a smooth spinor intertwiner $M$, the compressed forms $M^\ast f_\delta M$ and $M^\ast h M$ commute with the spin action, and are scalar by irreducibility. To average the distinguished ninth coordinate in the second contraction, compress the nine multiplication operators $f_\delta^2\sum_a X_{ai}^2$. Each commutes with the lifts of even sign flips after compression. The exact gamma-word commutant for those lifts is scalar. Rotations permute the spatial axes, so the nine scalars agree. Consequently the $X_{a9}^2$ contraction can be replaced by $r^2/9$ in each compressed expectation. Inserting the coefficient gives $(61-45/9)/1568=1/28$.

Choose a unit spinor weight vector with $J$-weight $2$, and let $\eta$ be its image. Write $a=\norm{B_\delta\eta}^2$ and $q=h[\eta]$. The inequality $f_\delta^2r^2\le f_\delta$ gives
\[
 \sum_\alpha\norm{f_\delta\tau_\alpha\eta}^2\le a/28,\qquad
 \sum_\alpha\norm{f_\delta\Xi_\alpha\eta}^2\le 8a .
\]
Taking the expectation in \eqref{eq:anticomm}, using source symmetry and Cauchy--Schwarz on the two sums, yields
\[
 2a\le 8\sqrt{qa/28}+2\sqrt{8/28}\,a .
\]
For $a>0$, divide by $2\sqrt a$, move the last term to the left and square. The case $a=0$ is immediate. The coefficient is
\begin{equation}
 \norm{B_\delta\eta}^2\le \kappa h[\eta],\qquad
 \kappa=\frac{4}{(\sqrt7-\sqrt2)^2}=\frac{36+8\sqrt{14}}{25}.
 \label{eq:hardy}
\end{equation}
Scalarity of the compressed forms extends this from the chosen weight to every spinor column. The graph-norm approximation from Chapter~\ref{ch:spin} and boundedness of $B_\delta$ extend it to form-domain intertwiners. The argument uses the actual covariance, source anticommutator and finite contractions listed above. Their combined source Hardy implication retains its analytical status.
\evidence{\src{research/computational/python-first-20261124/bfss-physical-implication/README.md}{End-to-end analytical source audit}, section 6. The source anticommutator and contraction receipts there include G14, G15 and X0.}

\section{Resolvent proof of the projection estimate}
\begin{theorem}[Conditional fidelity estimate]\label{thm:fidelity}
Assume the common square, the physical $\Spin(9)$ graph action of
Chapter~\ref{ch:spin}, and \eqref{eq:hardy} for physical spinor intertwiners. If $\psi$ is a unit physical spatial singlet in the weighted charge domain, then for each such $\delta>0$,
\[
 \norm{P_{\ker H}\psi}^2\ge 1-\kappa q_\delta[\psi].
\]
\end{theorem}
\begin{proof}
Let $t>0$, $\chi_\alpha=A_\alpha\psi$, and
$\eta_\alpha=(H+t)^{-1}\chi_\alpha$. Equation~\eqref{eq:chi} and resolvent covariance make $\alpha\mapsto\eta_\alpha$ a form-domain intertwiner. Write
\[
 a_\alpha=\ip{\chi_\alpha}{\eta_\alpha}
       =h[\eta_\alpha]+t\norm{\eta_\alpha}^2\ge0.
\]
The identity $\chi_\alpha=B_\delta w_\delta\chi_\alpha$, self-adjointness of the multiplier, and \eqref{eq:hardy} give
\[
 a_\alpha\le\norm{w_\delta\chi_\alpha}\norm{B_\delta\eta_\alpha}
       \le\norm{w_\delta\chi_\alpha}\sqrt{\kappa a_\alpha}.
\]
It follows, including the zero case, that
$a_\alpha\le\kappa\norm{w_\delta\chi_\alpha}^{2}$.

By $A_\alpha^2=H$, the spectral calculus for the self-adjoint $A_\alpha$ gives
\[
 a_\alpha=\int_{[0,\infty)}\frac{\lambda}{\lambda+t}\,d\mu_\psi^H(\lambda).
\]
Average over the sixteen labels and hold $\delta$ fixed while $t\downarrow0$. The integrand increases to $\mathbf1_{(0,\infty)}$, so the left side tends to $1-F(\psi)$. The right side is $\kappa q_\delta[\psi]$.
\end{proof}

If the trial lies in the weighted domain for all positive regulators and the same $\kappa$ works, each inequality bounds the same scalar $F(\psi)$. Taking the supremum of the scalar lower bounds gives $F(\psi)\ge 1-\kappa q_{\rm aux}[\psi]$. The fixed-regulator examples in Chapter~\ref{ch:certificate} already give conditional positive margins.

\chapter{Constructing and Checking an Exact Trial Witness}\label{ch:witness}
\chaplead{A finite variational calculation supplies a rational energy only after its rows and entries have been identified with the literal original charge form.}

\section{The packed seed and polynomial cores}
Let $J_{\rm pack}$ denote the isometry from packed occupation coefficients to the original Fock carrier. The packed formalization transports each original Majorana action on arbitrary coefficient vectors. It provides the interface through which exact finite occupation computations can be interpreted as statements about original Fock vectors. The trial seed is the normalized vector
\[
 \xi=J_{\rm pack}(S_{\mathbb Z})/\sqrt{8072064},
\]
with the stored integer coefficient vector $S_{\mathbb Z}$. Its cancellation under the color and spatial generators is a finite source condition needed for the singlet argument. The archive records exact sparse cancellation and selective independent replay.

Use three cyclically indexed color rows and define finite operators and polynomials by
\[
 K_{ai}=i\sum_{\alpha,\beta}(\gamma_i)_{\alpha\beta}
            t_{\alpha,a+1}t_{\beta,a+2},\qquad
 Z(X)=\sum_{a,i}X_{ai}K_{ai},\qquad Y(X)=Z(X)/\sqrt2.
\]
The row indices on the Majoranas are read cyclically. If $M_{dc}=\sum_iX_{di}K_{ci}$, one cubic core is defined by
\[
 2\sqrt2\,\Phi(X)=2\bigg(160sI-Z(X)^2
                  -\sum_{d<c}(M_{dc}-M_{cd})^2\bigg)Z(X)\xi .
\]
The second cubic, denoted $\Phi_3$, uses an ordered contraction of $X^{\otimes3}$ with three $K$ factors. Write $U$ for the sum over the six permutations of the color slots of $X^{\otimes3}$ and $D_{\rm c},D_{\rm s}$ for the trace-insertion sums in the color and spatial slots. Set $L=D_{\rm s}(5U-D_{\rm c}U)$ and define
\[
 2\sqrt2\,\Phi_3(X)
  =\frac1{330}\sum_{a,b,c,i,j,k}
     L_{abcijk}(X)K_{ai}K_{bj}K_{ck}\xi .
\]
The order of the $K$ operators is part of the definition. The literal tensor construction gives a trial using the original coefficients and does not require an identification with a Casimir projector from an independent construction.

The twenty-one cores are $Y^p\xi$ for $p=0,\ldots,8$, and $Y^p\Phi,Y^p\Phi_3$ for $p=0,\ldots,5$, with the latter families interleaved by degree in the archive. The source covariances of $K$, the seed and the tensor contractions imply equivariance of these cores under the stated actions. Multiplication by gauge- and rotation-invariant scalar factors preserves this law.

\section{The 104 polynomial-Gaussian rows}
Each row multiplies a core by a monomial $s^k u^\ell v^m$ and a sphere-averaged Gaussian. In physical width variables the latter is
\[
 f_{a,b}(X)=\int_{S^2}
  \exp\!\left(-\tfrac a2 s+\tfrac{a-b}{2} n^{\mathsf T}G n\right)d\sigma(n),
\]
with $\sigma(S^2)=1$. This is an orientation average with a specified measure. The row coefficients are exact decimal rationals from the frozen trial specification. In the source expression a row of total polynomial degree $d$ has the dilation factor $(29/20)^{d/2}$. An overall common dilation factor cancels when the full vector is normalized. The absolute sizes in the term-by-term factors still affect the quotient.

Positive widths, with smallest width $29/20$, give an upper bound by a polynomial times $\exp(-29r^2/40)$ for the trial and its differentiated terms. The unnormalized vector is a finite sum of these rows. Positivity of its squared norm is a finite test on its Gram matrix. Its normalized version is the $\psi$ used in Chapter~\ref{ch:certificate}.

To see membership in the weighted domain under the covariance premises, multiply the trial by smooth radial cutoffs. Their supports grow to exhaust configuration space. Both the original weighted outputs and the inputs converge by the Gaussian bounds. For the charge of the cutoff, the product rule adds a first-order symbol term. On its annulus $\norm{\nabla\chi_R}_\infty\le C/R$, so $w_\delta\norm{\nabla\chi_R}$ is bounded there for fixed $\delta$. The tail of the trial then makes this additional term go to zero in $L^2$. The cutoffs are genuine invariant smooth-core vectors when the gauge equivariance holds. Closure of the original graph and its weighted inverse-image description give the required membership.

\evidence{Packed carrier \run{45633634783}.
\src{research/computational/python-first-20261124/bfss-physical-implication/README.md}{Literal trial formulas, seed provenance, and cutoff argument}.}

\section{From finite contractions to global polynomials}
The computational inputs are homogeneous scalar or matrix-valued polynomials in $X$. Their degrees follow by differentiating the polynomial cores or inserting the specified linear and quadratic fields. Hermitian products use their real parts for the real row coefficients. Scalar contractions are invariant under $\SO(3)\times\SO(9)$. The two matrix families transform by color conjugation and are spatially invariant.

Every real $3\times9$ matrix can be brought to rectangular diagonal form by rotations in both factors. Orientation changes in singular-value decomposition can be compensated using sign choices and unused spatial axes. A sign change of any one diagonal entry is a spatial rotation when accompanied by a sign change in one unused direction. Thus diagonal restrictions are polynomials in the squared singular values $g_i$. An even simultaneous color and spatial sign flip fixes a rectangular diagonal matrix, and forces the off-diagonal entries of an equivariant color matrix to vanish there. Signed permutations give the corresponding permutation symmetry of the restrictions.

A symmetric polynomial in $g_1,g_2,g_3$ is a polynomial in $(s,u,v)$. The $i$th diagonal component of a matrix family is symmetric in the other two variables, so it is a polynomial in $g_i$, $s-g_i$ and $u-sg_i+g_i^2$. Reduce its powers by
\[
 g_i^3=sg_i^2-ug_i+v .
\]
Covariance shows that the same polynomial matrix is obtained globally from a sum of $I,G,G^2$ with symmetric polynomial coefficients. This argument includes repeated or zero singular values. The degree of each contraction now gives a finite monomial basis.

If the raw values of a contraction are correct at the fitted configurations, a full-column-rank design proves the coefficients in that basis. A nonzero determinant minor modulo a prime supplies exact rank over $\mathbb Q$. Held-out and independent direct values challenge the inputs and implementation. They have evidentiary value even when the rank already implies uniqueness of the coefficients.

The archive traces 81 original raw receipts through 79 full-rank designs, with 58,458 exact comparisons of recorded scalar values. That audit checks the coefficient lineage and individual held-out entries under the raw values. An independent source-exact replay of all raw contractions needed for a witness remains open. The main reduced target uses 52 rows and 12 active cores, with a 56-row fallback. Its existing rectangular-configuration derivative tests supply partial direct coverage. A full certificate must reach each nonzero fitted family with its required inputs or a proved structural relation.

\section{Weighted form and independent Gaussian integration}
In the source derivative expansion write $Q_\alpha=-iT_{\alpha d}\partial_d/\sqrt2+B_\alpha$, where $d$ indexes a spatial-color coordinate and $T_{\alpha d}$ is the linear combination of normalized Majoranas. Repeated $d$ in that expression is summed. Let $A_{de}=\sum_\alpha T_{\alpha d}T_{\alpha e}$ and $C_d=\sum_\alpha T_{\alpha d}B_\alpha$. For $W=r^2+\delta^2$ the weighted form is obtained by integrating the charge norm:
\begin{align*}
 q_\delta[\psi]={}&\int W\left(\tfrac12\norm{\nabla\psi}^2+
                        \Re\ip{\psi}{(u+Y)\psi}\right)dX\\
 &-\tfrac1{32}\Re\int\sum_{d\ne e}(\partial_d W)
                      \ip{\psi}{A_{de}\partial_e\psi}\,dX\\
 &-\Re\frac{i}{16\sqrt2}\int\sum_d(\partial_d W)
                   \ip{\psi}{C_d\psi}\,dX .
\end{align*}
The finite source expansion gives the coefficient tests for this formula. Symmetry of second derivatives cancels the antisymmetric off-diagonal kinetic contractions. Integration by parts on the mixed terms gives the original Yukawa field and the derivative of $W$. The two derivative-of-weight corrections are essential. Gaussian decay supplies the boundary control for these integrations.

The product of two oriented Gaussian profiles is $\exp(-\tr(KG))$ for positive $3\times3$ precision matrix $K$. Its density normalized in $X$ has nine independent column vectors with covariance $K^{-1}/2$. For symmetric test matrix $L$,
\[
 \mathbb E \exp(\tr(LG))=\det(I-K^{-1}L)^{-9/2}
\]
near $L=0$. Differentiating at zero gives Gram moments. Because the integrable monomials reduce to symmetric invariants and matrix insertions, one can generate the needed moments with exact rational arithmetic. After integrating the Gaussian in $X$, relative orientation reduces to a scalar integral in the cosine of the angle between the two color axes. The independent audit uses the integrated orientation primitive, the common positive scale of all the matrices, and its own moment recurrence from the determinant identity.

The independent Gaussian audit reproduced 203 moment polynomials through degree 11 and all 16,380 upper-triangle matrix entries of the three 104-row forms. These entries use the archived contraction polynomials. The rational endpoint computed from a coefficient vector $c$ has the schematic form
\[
 h[\psi]=\frac{c^{\mathsf T}M_Hc}{c^{\mathsf T}M_0c},
 \qquad
 q_{\rm aux}[\psi]=\frac{c^{\mathsf T}M_{\rm aux}c}{c^{\mathsf T}M_0c}.
\]
A common positive matrix normalization cancels. The correspondence between the raw occupation values and the literal contractions is still a load-bearing finite premise. The monograph keeps the verified integration and this remaining replay task visible in the final theorem.
\evidence{\src{research/computational/python-first-20261124/bfss-independent-integrals/README.md}{Independent Gaussian moment and matrix audit}.
\src{research/computational/python-first-20261124/bfss-physical-implication/NAG4_INTERVAL_DIRECT_CLOSED_PAIRING_AND_WITNESS_COVERAGE_2026_11_27.md}{Reduced witness coverage}.}

\chapter{The Conditional Existence Certificate}\label{ch:certificate}
\chaplead{The operator and computation meet in a statement about the kernel of a single Hamiltonian.}

\section{Exact statement and proof}
\begin{theorem}[Conditional zero-energy existence]\label{thm:exist}
Let $H$ be the original paired relative $\SU(2)$ Hamiltonian on $\E$ and let $\psi$ be the normalized 104-function trial of Chapter~\ref{ch:witness}. Suppose the following statements hold on the specified source.
\begin{enumerate}
 \item The natural squares of the self-adjoint individual charges equal $H$ on their exact domains.
 \item The physical $\Spin(9)$ vector and graph action gives the intertwining rule \eqref{eq:chi}, preserves $h$, and makes $\psi$ a singlet. The trial belongs to the weighted charge domain.
 \item The original physical spinor Hardy estimate \eqref{eq:hardy} holds at $\delta=1/2$.
 \item The literal source contractions and row evaluation identify the archived normalized matrix quotients with the two forms on $\psi$. They give
 \[
       q_{\rm aux}[\psi]<\frac{327658}{10^6},\qquad
       h[\psi]<\frac{36704}{10^6}.
 \]
\end{enumerate}
Then $\norm{P_{\ker H}\psi}^2>11.075824\%$. In particular there exists a norm-one vector in $\dom H$ of energy zero.
\end{theorem}

\begin{proof}
By Theorem~\ref{thm:fidelity},
\[
 F(\psi)\ge 1-\kappa\big(q_{\rm aux}[\psi]+\tfrac14 h[\psi]\big).
\]
Since $\sqrt{14}<15/4$, $\kappa<66/25$. The strict rational upper bounds give
\[
 1-\frac{66}{25}\frac{327658+36704/4}{10^6}
    =\frac{2768956}{25000000}=0.11075824 .
\]
The projection is therefore nonzero and its normalization gives the asserted vector.
\end{proof}

The same construction at unit regulator gives
$q_1[\psi]<3/8$ and $F(\psi)>1/100$ by the simple rational comparison in Chapter~\ref{ch:question}. Both estimates use the same trial, original $H$ and original kernel. The larger conditional bound uses the smaller positive regulator. The published preprint of Michel, Scheinpflug and Yin gives its own auxiliary-operator argument and variational result \cite{msy}. The present theorem uses the source-defined witness and the internal exact comparison. The preprint's numerical endpoint is not a premise of the theorem.

\section{Which evidence would settle each premise?}
\begin{description}
 \item[Common square.] Finish the selected native output and Hilbert-sum pairing assembly, then prove equality of weak-representative graphs. A consumer needs to exercise both the natural square domain and the action.
 \item[Covariance and admissibility.] Trace the original global wavefunction action from its exact infinitesimal conventions through the physical charge graph. Check each trial covariant, seed, and weighted cutoff against those conventions.
 \item[Hardy estimate.] Bind the exact anticommutators and contractions to the same $\Spin(9)$ representation used in the resolvent argument, including the compression that recovers the factor $1/9$ and the graph-norm intertwiner approximation.
 \item[Finite witness.] Independently replay a complete set of required raw source contractions, or provide exact certificates of the needed values. Use the existing rank and integration audits to close the remaining source-to-quotient link. Selected fixture agreement does not determine an entire family without the rank or a symbolic identity.
\end{description}

The ledger marks the source existence claim open. Strict closure results, ordinary proofs, and finite checks carry different evidentiary roles in this statement. Verifying the remaining source premises could yield a zero-mode result for the relative model. Questions about the multiplicity of the kernel, a gap above it and the interpretation of the model are separate.
\evidence{\src{research/computational/python-first-20261124/bfss-physical-implication/FORWARD_ANALYTICAL_RESEARCH_2026_11_25.md}{Exact positive-regulator endpoints}.
\src{research/methodology/SCIENTIFIC_KNOWLEDGE_ATLAS_PHYSICAL_EXISTENCE_PILOT.md}{Physical existence evidence ledger}.}

\part{Independent routes and physical comparisons}
\chapter{Localization, Escape, and Other Spectral Questions}\label{ch:local}
\chaplead{Localization can detect a zero mode by producing a tight small-charge sequence. A coercive sector estimate addresses a different spectral question.}

\section{First-moment tightness and the closed graph}
Consider original physical smooth-core vectors $f_n$ with $\norm{f_n}=1$, $Tf_n\to0$ in the Hilbert sum and
\[
 \sup_n\int r\norm{f_n(X)}^2\,dX<\infty .
\]
The last condition bounds the $L^2$ mass outside radius $R$ by $C/R$. The original kinetic identity \eqref{eq:energy}, nonnegativity of $u(X)$ and the linear bound on $Y$ then give a uniform gradient estimate. Applying local Rellich compactness to smooth cutoffs, and using the uniform tail bound, gives a strongly convergent subsequence in ambient $L^2$. The limit has norm one and lies in $\E$. Closedness of $T$ gives $Tf=0$, so $f\in\ker H$. This presents the conventional reasoning behind the accepted conditional criterion.

Exactly the same tail calculation explains why local compactness by itself does not give compactness over the full space. Translation or moving-annulus examples can have strong local convergence to zero and unit norm globally. The BFSS valley gives a source-motivated region to test for escape, but a formal valley packet with bounded form energy is not an operator-domain Weyl sequence.
\evidence{GM5E6 R2 qualifies the tight-core zero-mode implication. \run{45191316189}. The underlying first-moment compactness needs the specified derivative and closed-domain hypotheses.}

\section{A fixed-ball alternative}
A normalized small-charge sequence need only retain some uniformly positive mass in one fixed ball if one can localize with a suitable controlled IMS partition. For the original gauge-invariant smooth core, the cutoff product rule controls how $T$ changes under multiplication. A fixed smooth gauge-invariant partition $\{\chi_j\}$ with $\sum_j\chi_j^2=1$ cancels the mixed terms because $\sum_j\chi_j\nabla\chi_j=0$. The remaining term is bounded by the derivative weights of the partition and the original Clifford principal symbol.

GM5E10 R2 gives the precise conditional conclusion. On the original invariant core let the sequence have input norms at most one, native charge tending strongly to zero, and positive mass uniformly retained in a fixed closed ball. Given the finite smooth gauge-invariant IMS partition with its required support and gradient conditions, the theorem yields a nonzero zero mode on the natural Hamiltonian domain. Local compactness supplies persistent overlap after an appropriate cutoff, and the closed graph carries the low-charge conclusion. A sequence entirely supported in moving disjoint annuli cannot satisfy the fixed-ball hypothesis.

Explicitly, the inner member $\chi_{j_0}$ equals one on $r\le R$ and vanishes for $r>S$, with $S\ge0$. The finite family is smooth and gauge invariant, with $\sum_j\chi_j^2=1$ and $\norm{\nabla\chi_j}^2\le D_*$ for all members and all configurations. For input norm at most one, the archived cutoff estimate bounds the squared \emph{unnormalized} charge column of the inner member by $\norm{Tf_n}^2+8\,|\{j\}|D_*$. Its radial moment is at most $S$, and its squared norm is bounded below by the mass in the fixed ball. The localized sequence serves as a compact test of overlap with the uncut sequence. Its charge is bounded; the charge that tends to zero is the uncut charge.

In ordinary Euclidean coordinates a pair of partition functions can be constructed from a smooth angle $\varphi(r^2)$ that equals zero through $R^2$ and $\pi/2$ from $S^2$ onward, with $0\le R<S$. Take the cosine as the inner member and the sine as the other member. They are smooth and gauge invariant, their squares sum to one, and their gradients are bounded since their derivatives are supported in a fixed annulus. A native Lean instantiation of this construction and a source trial sequence meeting the probability and charge hypotheses remain separate tasks.
\evidence{\src{research/integration/BFSS_GM5E10_R2_ORIGINAL_FIXED_BALL_ZERO_MODE_GREEN_ACCEPTANCE_2026_11_23.md}{GM5E10 R2 qualification and original artifact}. \run{45236470826}. This succeeds the fixed-ball proof obligation listed in Edition 0.9.1.}

\section{Conditional high-type confinement}
For a reducing physical sector $E_s$, an estimate of the form
\begin{equation}
 h[f]+b\norm{f}^2\ge
      a\norm{\nabla f}^2+c\int r\norm{f(X)}^2\,dX,\qquad a,c>0,
 \label{eq:coercivity}
\end{equation}
on the actual closed form domain would make its form embedding compact by the argument above. Functional calculus then makes the sector resolvent compact. For example, the resolvent maps bounded sets to a bounded set in the form norm, and the right side of \eqref{eq:coercivity} makes that set precompact in $L^2$. A positive eigenpair on the full physical space would also require a nontrivial reducing sector with suitable positive spectral content. The corresponding $\Spin(9)$ vector branching, transverse oscillator estimate, fermionic constant and domain transport remain distinct steps.

The transverse lower bound begins at a fixed nonzero first row $X^1$ with length $\rho$ and direction $a$. Decompose the other rows as $t_2a+z_2$ and $t_3a+z_3$ with $z_i\perp a$. The wedge potential is at least $\rho^2(\norm{z_2}^2+\norm{z_3}^2)$. At fixed first row the change of coordinates is orthogonal on the other two rows, and their two perpendicular components have sixteen real coordinates. For $0<\beta<1/2$, reserve $\beta$ of the kinetic coefficient for $-\beta\Delta_z+\rho^2\norm z^2$ and retain $1/2-\beta$ in all coordinates. The auxiliary oscillator levels are $2\sqrt\beta\,\rho(n+8)$. This slice calculation differentiates only the two remaining rows. A specified physical sector would have to exclude enough low transverse levels for the resulting slope to dominate the unnormalized linear fermion bound. The scalar oscillator spectrum alone gives no such exclusion.

The angular average used in that proposed estimate has a useful exact control. If $\sigma$ is normalized measure on $S^2$ and $G=XX^{\mathsf T}$, then
\[
 \int_{S^2}\norm{v^{\mathsf T}X}\,d\sigma(v)\ge \tfrac12\norm{X}.
\]
For $X\ne0$ diagonalize $G$ and write its eigenvalues as $\norm{X}^2p_i$ with $\sum p_i=1$. Concavity gives $\sqrt{\sum_i p_i v_i^2}\ge\sum_i p_i|v_i|$, and the integral of $|v_i|$ equals $1/2$. A rank-one matrix attains equality. On a physical norm density invariant under full adjoint color rotations the integrated directional energy is independent of $v$, so the average applies to any fixed direction under the source binding. The coefficient $1/2$ is a real constraint on any oscillator margin. Recovering the correct physical Fock-vector sector and showing that the estimate survives closure remain essential.

The conditional zero-mode criteria and conditional compactness of a high-spin sector address different endpoints. An external numerical construction of ground-state or excited-state wavefunctions gives another comparison point \cite{bajajlin}. Agreement of observables or energies must be tied to the specific operator and domain to support a source theorem.

\chapter{Relational Topology and Particle Structure}\label{ch:topology}
\chaplead{A geometric structure becomes relevant to the model through a specified map, its gauge behavior, and its relation to dynamics.}

\section{Incidence and the classical color orbit}
View the columns of $X$ as vectors $X_i\in\R^3$ in color space. The full adjoint rotation action preserves $X_i\cdot X_j$ and the signed triples $\det(X_i,X_j,X_k)$. The signed triples use orientation preservation. The source incidence results construct the relevant $\SO(3)$ adjoint action and its invariant data.

There is an ordinary finite-dimensional converse for these explicit arrays. Suppose two configurations have the same column Gram matrix. Every linear relation among one set of columns is also a relation among the other, as can be checked by squaring the norm of the corresponding linear combination. Sending the columns to their counterparts defines an isometry of their spans. Extend it orthogonally to $\R^3$. If the span has dimension below three, the extension can be chosen with determinant one by reversing an unused direction if necessary. If it has dimension three, take a linearly independent triple. Equality of the signed determinants forces the isometry to have determinant one. Thus the arrays distinguish classical color orbits of the explicit $\SO(3)$ action. A source implementation of the converse and a map to quantum physical states have separate proof requirements.

For a rank-one configuration write $X=c\,n^{\mathsf T}$ with $c\ne0$ and $\norm n=1$. There is a color rotation carrying $c$ to $-c$. The coarse direction after this identification is the unoriented line $[n]\in\mathbb{RP}^8$. This direction omits radial data and stabilizers of the original action.

\section{Gamma eigenspaces and projective descent}
On the oriented sphere define $\Gamma(n)=\sum_i n_i\gamma_i$ and $P_\pm(n)=(I\pm\Gamma(n))/2$. The Clifford identity gives $\Gamma(n)^2=I$ and its eigenspaces have rank eight. For example, a unit vector perpendicular to $n$ gives a gamma matrix that anticommutes with $\Gamma(n)$ and exchanges its two eigenspaces. The antipodal map also exchanges them:
\[
 P_+(-n)=P_-(n).
\]
A literal color arrow acts on the color factor and trivially on the sixteen-spinor labels. It commutes with a fixed spinor projection. If it intertwined the individually labelled $P_+(n)$ with the individually labelled $P_+(-n)$ under the rank-one identification, it would force $P_+(n)=P_-(n)$, which contradicts their complementary nonzero ranges. The obstruction concerns that labelled descent.

There is a compatible twisted field. If $\ell$ denotes the tautological line over $\mathbb{RP}^8$, a vector $z\in\ell_{[n]}$ determines the map $z\mapsto\Gamma(z)$. In a local choice of unit generator the sign of the matrix changes with the sign of the generator. This is precisely the transition law of the line factor. The unordered pair of spectral projections is also compatible with reversal.

An additional conditional test uses characteristic classes. If the source $+$ bundle over $S^8$ has $w_8\ne0$, it cannot be the pullback of an ordinary bundle from $\mathbb{RP}^8$. The double covering pulls top mod-two cohomology back by degree $2=0$ modulo two, and naturality would force the top class to vanish. Lower-degree positive cohomology of $S^8$ is zero, so the color-triple sum retains this $w_8$ by the Whitney formula. The source clutching comparison required to establish the premise belongs to the archived topological analysis.

\section{Qualified relational topology and its interpretation}
The classical octonionic Hopf fibration has spheres $S^{15}$ and $S^8$ with fiber $S^7$ \cite{gwz}. The original Gamma9 quadratic norm identity and the projective constructions use exact source data. D69 qualifies the homeomorphism from a particular original source-$Q$ circle orbit quotient to its canonical complex projectivization. This statement identifies a quotient with a specific target. A Hopf or Chern invariant, a physical wavefunction carrier map, and dynamical protection of a texture require their own constructions. A map from a three-sphere into $S^8$ alone supplies no homotopy protection, since $\pi_3(S^8)=0$.

Related proposals for nested Hopf geometry or protected persistence can use the same descent test. They must specify the physical carrier, allowed variations, a conserved or topological quantity, and a mechanism connecting it to the dynamics. The source quotient is a precise mathematical comparison point.
\evidence{D69 R1 \run{45048282121}.
\src{research/topology/RHOPF_P0_RELATIONAL_OCTONIONIC_QUATERNIONIC_HOPFION_AUDIT_2026_11_18.md}{Gamma9 and homotopy audit}.
\src{research/unification/SPIN9_HOPF_BUNDLE_PHYSICAL_VALLEY_RP8_OBSTRUCTION_2026_11_10.md}{Physical valley direction and antipodal obstruction}.}

\section{Color labels, quadratic invariants and chirality}
In the original real Majorana label space let $C=\R^3$ carry the adjoint color action and $S=\R^{16}$ the spatial spinor labels. An invariant rank-one operator on $C$ would give an invariant line, which irreducibility of the standard $\SO(3)$ action forbids. Thus three separately gauge-invariant rank-one color projectors cannot serve as the proposed family labels.

Quadratic invariants do exist. The real decomposition
\[
 \Lambda^2(C\otimes S)
  \cong(\Lambda^2 C\otimes\operatorname{Sym}^2S)
      \oplus(\operatorname{Sym}^2C\otimes\Lambda^2S)
\]
shows that the gauge-invariant part is $\Lambda^2 S$. Here $\Lambda^2 C\cong C$ has no invariant vector, and the only invariant in $\operatorname{Sym}^2 C$ is the metric up to scale. Color contraction therefore gives $120$ real skew-quadratic coefficient directions. A physical Fock observable associated with any such direction requires a quantization, a physical domain, and its own dynamical calculation.

Both complex chiral spinors of $\Spin(10)$ restrict to the same complex $\Spin(9)$ spinor. This abstract restriction loses the distinction between the two chiralities. The rotational spinor label therefore supplies no four-dimensional chirality through dimension matching alone. The gauge projection, the meaning of Lorentz and internal gauge actions, and the large-scale observables all belong to the proposed physical map.
\evidence{\src{research/unification/BFSS_THREE_BY_SIXTEEN_SPIN9_SPIN10_FAMILY_FEASIBILITY_2026_11_10.md}{Family feasibility audit}.
\src{research/methodology/BFSS_UNIFIED_QUALIFIED_INCIDENCE_SPIN9_MERGE_CHECKPOINT_2026_11_15.md}{Qualified incidence integration}.}

\chapter{Standard Model Benchmarks and Tests of Emergence}\label{ch:comparison}
\chaplead{A comparison with nature requires a source-defined map to an observable and a prediction that can discriminate among admissible alternatives.}

\section{The representation benchmark}
For a four-dimensional chiral matter claim, one benchmark is a single Standard Model family in left-handed notation. In the convention $Q_{\mathrm{em}}=T_3+Y$, take
\[
 (3,2)_{1/6},\quad(\overline3,1)_{-2/3},\quad
 (\overline3,1)_{1/3},\quad(1,2)_{-1/2},\quad(1,1)_1 .
\]
A gauge singlet with $Y=0$ would have no effect on the anomaly arithmetic. The representations, chirality and multiplicities are inputs. With common Dynkin-index normalizations suppressed, the two nonabelian mixed hypercharge sums are
\[
 2(1/6)-2/3+1/3=0,\qquad 3(1/6)-1/2=0 .
\]
The abelian sums $\sum dY$ and $\sum dY^3$ over representation dimensions $d$ vanish as well. For example the cubic sum is
\[
 6(1/6)^3+3(-2/3)^3+3(1/3)^3+2(-1/2)^3+1=0 .
\]
The color cubic sum is $2-1-1=0$. These exact cancellations are tests of specified input content. A candidate theory still needs to derive its content, its multiplicities across families, and its relation to measured observables. The electroweak, Higgs, flavor and oscillation reviews provide empirical and phenomenological reference points \cite{pdg}.

The project's empirical ledger records provenance, conventions and uncertainties. The ledger qualification concerns the curated records. Experimental outcomes, derived model values, covariance matrices and a statistical likelihood play distinct roles. This edition makes no original candidate-derived Standard Model prediction.

\section{A finite test for distinguishability}
Structural Counterfactual Capacity (SCC) asks which differences between admissible candidate inputs survive a specified observable map. In a finite setting take two admissible candidate families and a set of valid measurement protocols. With exact categorical predictions $p,q$ for a shared protocol, the algebraic distinction is
\[
 d_{\mathrm{TV}}(p,q)=\tfrac12\sum_i|p_i-q_i|>0 .
\]
A collection of protocols distinguishes the two families uniformly when every cross-family pair has at least one validated separating protocol. Missing prediction entries cannot provide a witness. This exact comparison of predictions leaves the size of a statistically resolvable difference to the sampling model and the empirical analysis.

A simple negative control in neutrino phenomenology is the absolute offset. An observable depending only on $m_2^2-m_1^2$ is unchanged by a common admissible shift of the two squared masses. Its exact value therefore leaves that offset undetermined. A BFSS-to-neutrino proposal must specify the relevant measured quantity, identify its source states, and give the convention and physical units. An internal analogue is the invisibility of a symmetric trace-free coefficient perturbation under a quadratic CAR contraction. Both examples locate an ambiguity in the observable map.

\evidence{\src{research/standard-model/SM_D0_R1_QUALIFIED_EMPIRICAL_LEDGER_ACCEPTANCE_2026_11_18.md}{SM-D0 empirical ledger}.
\src{research/standard-model/SMD1_SCC3B1_R1_QUALIFIED_ORTHOGONAL_EXACT_JUDGMENT_METHODOLOGY_2026_11_18.md}{SCC3B1 exact comparator}.}

\section{An admissible test of relational state identity}
The working Hypothesis of Everything asks whether persistent identities, symmetries and interactions may emerge from a common relational quantum dynamics. Its November review proposes a finite, falsifiable subquestion. Given physically permitted Hermitian witnesses $\mathcal A$ and a group $\mathcal U$ of permitted unitary transformations, let
\[
 W=\operatorname{span}_{\R}\{I,\,U^\ast A U:
                       A\in\mathcal A,\ U\in\mathcal U\}.
\]
Specify admissible states and physical equivalence independently. Two density matrices have the same accessible ideal expectations exactly when their difference is orthogonal to $W$ for the real Hilbert--Schmidt pairing. To identify physical equivalence by these records, prove both invariance on equivalent states and separation of all inequivalent states.

If all density matrices are admissible and $W$ is proper, take nonzero Hermitian $B\in W^\perp$. Since $I\in W$ its trace vanishes. For sufficiently small $\epsilon>0$ the two matrices
\[
  \rho_\pm=I/d\pm\epsilon B
\]
are positive of trace one and give the same records. One can use
$\epsilon\le 1/(d\norm B_{\op})$. This gives a concrete indistinguishability certificate, subject to admissibility of the constructed states. Exact full Clifford algebra reconstruction in Chapter~\ref{ch:source} gives a useful upper-access comparison. The actual physical witnesses may be limited by gauge symmetry, parity and the permitted dynamics.

For finitely many Hermitian continuous-control generators $H_j$, compute the reachable real span by closing the initial span under $A\mapsto i[H_j,A]$. Finite dimensionality makes this process terminate when rank stabilizes. Invariance under the derivatives implies invariance under the exponential series. Conversely, differentiating the unitary orbits within their finite-dimensional closed span shows that it contains the derivatives. The two spans therefore coincide. In the archived two-level example the initial witnesses are $\{I,\sigma_z\}$. With no controls the rank is two. Adding the $\sigma_x$ generator exposes $\sigma_y$ and gives rank three. Adding $\sigma_z$ as another generator exposes $\sigma_x$ and gives rank four. The example tests the method in a model with specified permissions.

\section{Persistence and physical comparison}
Accessible state separation and persistence in time are different mathematical claims. Persistence needs a conserved quantity, an invariant or reducing sector, or a quantified stability result under stated perturbations. In the source relative model the first task is to identify an admissible witness/control system with a justified relation to the original physical Hilbert space and its dynamics. A vector-level kernel result could inform that task, while a comparison to a four-dimensional gauge theory would still need an effective map and a prediction.

Any proposal for a GUT, nested fibrational dynamics or a broader unification may be tested by the same protocol. State the observable and admissible alternatives, derive the candidate prediction, and retain an exact negative control where the observable cannot separate them. This makes the scope of each unresolved physical question explicit.
\evidence{\src{research/unification/HYPOTHESIS_OF_EVERYTHING_STATUS_AND_DYNAMIC_WITNESS_TEST_2026_11_27.md}{Hypothesis of Everything evidence review and reproduced Pauli control}.
\src{research/quantum-information/BFSS_QIT_SCC_W0H1B_R2_GREEN_44164974919_ACCEPTANCE_2026_11_17.md}{Full Clifford commutant qualification}.}

\appendix
\chapter{Notation and Source Normalization}\label{app:notation}
\small
\begin{longtable}{@{}p{.23\textwidth}p{.72\textwidth}@{}}
\toprule Symbol & Meaning and convention\\
\midrule
\endhead
$X,X_i,X^a$ & Matrix of color rows and spatial columns, its $i$th column in $\R^3$, and its $a$th row in $\R^9$.\\
$r,G,s,u,v$ & Frobenius radius, color Gram matrix $XX^{\mathsf T}$, and its three elementary symmetric functions in degree order. The incidence argument uses instead the column Gram data $X^{\mathsf T}X$.\\
$\F,\Ltwo,\E$ & Complex Fock vector space, ambient square-integrable sections, and the closed gauge fixed space. $\operatorname{End}_\C\F$ is the separate Fock operator space.\\
$\Core,\PhysCore$ & Compactly supported smooth sections and their gauge-equivariant subspace. The latter embeds densely in $\E$.\\
$\gamma_i,\theta,t$ & Real gamma matrices on sixteen spin labels, original source Majoranas, and normalized Majoranas $t=\sqrt2\,\theta$.\\
$Q_\alpha,A_\alpha,T$ & Core charge, its native self-adjoint closure on $\E$, and the original closed Hilbert-sum column.\\
$h,H$ & $h(f,g)=\ip{Tf}{Tg}/16$ on $D(T)$, and the normalized adjoint square with natural domain. Inner products are conjugate-linear in the first argument.\\
$Y(X),K_{ai}$ & Original unnormalized Yukawa multiplier and the finite operators used to define it in the trial coordinates.\\
$w_\delta,B_\delta$ & Positive radial weight $\sqrt{r^2+\delta^2}$ and bounded reciprocal for $\delta>0$.\\
$q_\delta,q_{\rm aux}$ & Weighted charge energy and the scalar form $\norm{rT\psi}^2/16$ on the trial domain.\\
$t$ in $(H+t)^{-1}$ & Positive resolvent shift. Its limit in the projection argument is independent of the radial regulator.\\
$F(\psi)$ & Squared projection of a unit trial onto the original kernel.\\
\bottomrule
\end{longtable}
\normalsize

The normalizations used for the energy calculation are the Pauli trace convention $e_a=\sigma_a/\sqrt2$, the ordered-pair $1/2$ coefficient of the core charge, and the factor $1/16$ on the Hilbert-sum pairing. The ordinary energy has kinetic factor $1/2$. The regulator is expressed as $\delta^2$ inside the radial weight. A formula written using an additive parameter under the square root therefore needs the corresponding change of parameter.

\chapter{Evidence Index and Declaration Map}\label{app:evidence}
The original 59-record \src{contributions/bfss-kotler/internal-monograph/THEOREM_ATLAS.json}{historical theorem atlas} preserves the earlier acceptance history. The 22-record \src{research/methodology/SCIENTIFIC_KNOWLEDGE_ATLAS_PHYSICAL_EXISTENCE_PILOT.json}{physical existence atlas} links the new chapters to their scientific prerequisites. The table presents the latter records as of the edition cutoff. Strict rows refer to the cited tested source and consumer; the archive retains their exact status fields. The dependency links express mathematical use and may differ from the Lean import graph.

\small
\begin{longtable}{@{}p{.25\textwidth}p{.47\textwidth}p{.21\textwidth}@{}}
\toprule Record & Claim & Evidence\\
\midrule
\endhead
\texttt{PE-UGC4} & Literal restricted joint graph equals full graph physical-input slice & Strict\newline\run{45869599441} \\
\texttt{PE-CSQE0} & Equal original core energies for every spin label & Strict\newline\run{45505859621} \\
\texttt{PE-ICDA} & Exact componentwise original joint graph adapter & Strict\newline\run{45884276018} \\
\texttt{PE-ICDB} & Individual closed graph lifts to selected joint coordinate under equal core norms & Strict\newline\run{45886320206} \\
\texttt{PE-ICDC} & Actual paired individual and joint input domains equal; selected component graph lift & Strict\newline\run{45912386053} \\
\texttt{PE-LOC6C} & Physical individual closed graph iff invariant-core weak pairing & Strict\newline\run{46041935392} \\
\texttt{PE-NAG0} & Actual native closed individual graph & Strict\newline\run{46047069230} \\
\texttt{PE-NAG1} & Native individual output physical; common dense domain & Strict\newline\run{46051319682} \\
\texttt{PE-NAG2} & Native individual self-adjointness on actual paired physical space & Strict\newline\run{46059907504} \\
\texttt{PE-NAG3} & Literal native smooth-core operator closure equals NAG2 operator & Strict\newline\run{46067320278} \\
\texttt{PE-ESA} & Explicit literal core essential self-adjointness & Strict\newline\run{46074553299} \\
\texttt{PE-CSQF0} & Complex core and closed joint graph pairing & Strict\newline\run{46082764637} \\
\texttt{PE-SQUARE} & Normalized common form and natural square equality & Analytical \\
\texttt{PE-SPIN} & Genuine Spin(9) Fock-vector and physical-domain covariance & Analytical \\
\texttt{PE-TRIAL} & Literal 104-function trial and weighted admissibility & Analytical \\
\texttt{PE-FINITE-LINEAGE} & Exact archived contraction-coefficient lineage and full-rank interpolation & Finite audit \\
\texttt{PE-RAW} & Independent source-exact raw contraction/seed certificate & Open \\
\texttt{PE-HARDY} & Positive-regulator Hardy estimate on physical spinor intertwiners & Analytical \\
\texttt{PE-FIDELITY} & Same-H fidelity inequality for all positive regulators & Analytical \\
\texttt{PE-EXIST} & Nonzero original physical zero-energy eigenstate & Open \\
\texttt{PE-NEG-ICD0} & Stopped monolithic ICD0 implementation & Stopped \\
\texttt{PE-NEG-GM2A} & Limit of the massless multiplier domain theorem & Scope limit \\
\bottomrule
\end{longtable}
\normalsize

\section{Source declaration map for the main proof}
The following names identify the source interfaces used in Chapters~\ref{ch:ham}--\ref{ch:square}. Each link is pinned to the edition snapshot; the linked qualification receipts specify their individual tested commits.

\begin{description}
\item[Original data.] The \src{lean/OAI/MathematicalPhysics/BFSS/Core.lean}{core file} defines the actual charged differential expression, polarization and core symmetry.
\item[Ambient joint graph.] ICD-A0 \nolinkurl{source_joint_closed_graph_eq} identifies the closure of the physical core range. ICD-B0 \nolinkurl{source_component_graph_iff} gives the selected-component lift. \src{lean/OAI/MathematicalPhysics/BFSS/Contrib/Kotler/BFSSICDA0OriginalJointGraphAdapter.lean}{ICD-A0 source} and \src{lean/OAI/MathematicalPhysics/BFSS/Contrib/Kotler/BFSSICDB0OriginalComponentGraphLifting.lean}{ICD-B0 source}.
\item[Physical graph and square.] AY gives the physical Hilbert-sum column, its graph equality and the membership equivalence under the coordinate map. BE proves the normalized Hamiltonian weak-form equivalence and the generic natural adjoint-square characterization. \src{lean/OAI/MathematicalPhysics/BFSS/Contrib/Kotler/OperatorDependencies/SU2BFSSPhysicalDomainL6H5E8AYOriginalPhysicalHilbertChargeGraphProbe.lean}{AY source} and \src{lean/OAI/MathematicalPhysics/BFSS/Contrib/Kotler/OperatorDependencies/SU2BFSSPhysicalDomainL6H5E8BEOriginalNormalizedPhysicalHamiltonianProbe.lean}{BE source}.
\item[Individual graph and domain.] NAG0 \nolinkurl{paired_individual_graph_iff}, NAG1 \nolinkurl{paired_native_domain_eq}, and NAG2 \nolinkurl{paired_native_selfAdjoint}. See \src{lean/OAI/MathematicalPhysics/BFSS/Contrib/Kotler/BFSSNAG0OriginalIndividualPhysicalOperator.lean}{NAG0}, \src{lean/OAI/MathematicalPhysics/BFSS/Contrib/Kotler/BFSSNAG1OriginalIndividualPhysicalDomain.lean}{NAG1} and \src{lean/OAI/MathematicalPhysics/BFSS/Contrib/Kotler/BFSSNAG2OriginalIndividualPhysicalAdjoint.lean}{NAG2}.
\item[Complex graph pairing.] CSQ-F0 \nolinkurl{paired_core_pairing}, \nolinkurl{pairing_on_closure}, \nolinkurl{joint_range_pairing}, and \nolinkurl{paired_closed_joint_pairing}. \src{lean/OAI/MathematicalPhysics/BFSS/Contrib/Kotler/BFSSCSQF0OriginalClosedJointPairing.lean}{CSQ-F0 source} and \src{lean/diagnostics/BFSSCSQF0OriginalClosedJointPairingIndependentAudit.lean}{Independent consumer}.
\end{description}

For the final analytical assembly, the normalized pairing uses \texttt{PiLp.inner\_apply} at exponent two from \href{https://github.com/leanprover-community/mathlib4/blob/d13f23b723b8a846827a245b89c10fc7d3f11612/Mathlib/Analysis/InnerProductSpace/PiL2.lean}{the pinned Mathlib source}. The original normalized Hamiltonian is separately indexed by E08-BF. The common-square conclusion itself is PE-SQUARE and remains analytical at the cutoff.

\chapter{Historical Lineage and Limits of Earlier Arguments}\label{app:lineage}
\section{How the current proof developed}
The October source work built the finite gamma and Fock carriers, identified the invariant physical core and related the undeformed energy to the literal core charge. S4B36 recovered the concrete core Pauli/Fock energy identity. The earlier H-03 comparison identifies historical and native paired operators under its source hypotheses. Its original qualification and the modified port must be tracked by their respective source records. The current normalized operator has its own exact integration evidence.

The November sequence recovered larger parts of the original gauge, adjoint and graph arguments in bounded interfaces. The weighted charge construction and shared kernel gave the fidelity route a precise target. The fixed-ball localization sequence advanced from GM5E9 to the qualified conditional GM5E10 criterion. NAG4 established essential self-adjointness of each literal core charge. CSQ-F0 replaced an unqualified closed-pairing step with a qualified two-point graph theorem. This chronology explains the order of investigation. It does not mean every earlier report is a logical premise of each later claim.

\section{Negative results and controls retained}
\begin{longtable}{@{}p{.3\textwidth}p{.65\textwidth}@{}}
\toprule Test or boundary & Consequence for reuse\\
\midrule
\endhead
CSQ-F0 R2 recursion limit & The first concrete core pairing specialization exceeded Lean recursion depth. R3 reused direct source symmetry and kept the topological lemma. The R2 failure concerns the implementation.\\
Monolithic ICD0 & The combined implementation stopped unqualified. The separate ICD-A0, ICD-B0 and ICD-C0 claims have their own receipts.\\
Massless multiplication & A maximal multiplier domain is a domain for the fermion field. The first-order weak charge graph needs distributional testing and graph approximation.\\
Bounded normalized Gram field & Its uniform bound applies to its normalized auxiliary expression. The original fermionic multiplication grows linearly with radius.\\
Moving annulus & Uniform norm or small form energy alone cannot give persistent probability in a fixed ball.\\
Finite witness fixture & Exact agreement on selected values tests those values. Global identification needs the admissible polynomial class, sufficient design rank and correct source inputs.\\
Three color projectors & Full adjoint rotations forbid three separately invariant rank-one color family labels. Quadratic color contractions remain possible.\\
Antipodal valley direction & Individually labelled gamma eigenspaces exchange under reversal. Twisted or unordered data have different descent conditions.\\
Full Clifford access & Algebraic separation of all finite states assumes access to all operators. Physical parity, gauge and dynamical constraints define a different set of permitted tests.\\
\bottomrule
\end{longtable}

\evidence{\src{research/computational/python-first-20261124/bfss-physical-implication/CSQF0_R2_NATIVE_RED_R3_DIRECT_CORE_PAIRING_2026_11_27.md}{CSQ-F0 R2 archive}.
\src{contributions/bfss-kotler/internal-monograph/appendices/negative_knowledge.tex}{Edition 0.9.1 negative-knowledge appendix}.}

\chapter{Edition Crosswalk and Reproduction}\label{app:crosswalk}
The rewrite groups the earlier accounts by their role in the existence argument. The table is a locator for prior chapter citations. Material with a historical qualification remains attached to the original receipt.

\small
\begin{longtable}{@{}p{.48\textwidth}p{.47\textwidth}@{}}
\toprule Edition 0.9.1 chapter & Edition 0.10.0 destination\\
\midrule
\endhead
1. Scope and physical context & Chapter~\ref{ch:question}.\\
2. Carriers and notation & Chapter~\ref{ch:source} and Appendix~\ref{app:notation}.\\
3. Clifford and Fock realization & Chapter~\ref{ch:source}, with packed transport in Chapter~\ref{ch:witness}.\\
4. Sixteen charges and smooth-core form & Chapters~\ref{ch:source} and \ref{ch:ham}.\\
5. Gauge symmetry and physical space & Chapters~\ref{ch:source}, \ref{ch:ham}, and \ref{ch:charges}.\\
6. Native physical Hamiltonian & Chapters~\ref{ch:ham} and \ref{ch:square}.\\
7. Physical massless multiplier & Chapter~\ref{ch:ham}.\\
8. Source Spin(9) and axis stabilizer & Chapters~\ref{ch:spin} and \ref{ch:topology}.\\
9. Fock-operator action and normalized Gram field & Chapters~\ref{ch:spin} and \ref{ch:ham}.\\
10. Gauge compatibility and sector reduction & Chapters~\ref{ch:spin} and \ref{ch:local}.\\
11. Original energy and fermionic growth & Chapters~\ref{ch:ham} and \ref{ch:local}.\\
12. First-moment compactness & Chapter~\ref{ch:local}.\\
13. Weighted charges and fidelity & Chapters~\ref{ch:fidelity}--\ref{ch:certificate}.\\
14. Conditional oscillator exclusion and spectrum & Chapter~\ref{ch:local}.\\
15. Hopf geometry and gauge descent & Chapter~\ref{ch:topology}.\\
16. Representation structure and families & Chapter~\ref{ch:topology}.\\
17. Standard Model and counterfactual tests & Chapter~\ref{ch:comparison}.\\
18. Conditions for a GUT or ToE & Chapter~\ref{ch:comparison}.\\
19. Proof recovery and negative results & Appendices~\ref{app:evidence} and \ref{app:lineage}.\\
20. The theorem frontier & Chapters~\ref{ch:certificate}, \ref{ch:local}, and \ref{ch:comparison}.\\
\bottomrule
\end{longtable}
\normalsize

\section{Reproducing the edition and its evidence}
The versioned manuscript directory contains \texttt{main.tex}, chapter sources, an evidence index, and build instructions. In a standard TeX Live installation run \texttt{latexmk -pdf main.tex}. Use the accompanying edition audit to check the include graph, bibliography and cross-references. The build record gives the PDF hash, page count and inspection scope. The manuscript build checks the exposition artifact. Exact Lean qualifications come from the archived source and independent consumer runs.

For a finite energy assertion, first locate the frozen seed, exact row coefficients, normalization, contraction receipts and matrix generation audit. Check the stated dependencies in the physical existence atlas. Reproducing the exact rational quotient from the archived matrices exercises a different step from recomputing the raw source contractions. Both steps retain their individual evidence ceilings.

For a source theorem, locate its declaration, original run and tested commit. Confirm source, consumer and workflow blobs, pinned environment and allowed axiom roots. Historical failures remain evidence about the attempted statements and proof architecture. The dated \src{contributions/bfss-kotler/internal-monograph/EDITION_0_10_0_REWRITE_PROSPECTUS_2026_11_27.md}{rewrite prospectus} gives the earlier editorial plan and the current source map that informed this edition.

\backmatter
\begin{thebibliography}{9}
\addcontentsline{toc}{chapter}{Bibliography}
\bibitem{bfss}
T. Banks, W. Fischler, S. H. Shenker, and L. Susskind.
\newblock M theory as a matrix model: A conjecture.
\newblock \emph{Physical Review D} \textbf{55} (1997), 5112--5128.
\newblock \href{https://doi.org/10.1103/PhysRevD.55.5112}{doi:10.1103/PhysRevD.55.5112}.

\bibitem{msy}
A. Michel, J. Scheinpflug, and X. Yin.
\newblock 99.9\ldots\% of the $\SU(2)$ BFSS ground state.
\newblock Preprint, November 2026.
\newblock \href{https://arxiv.org/abs/2611.09837}{arXiv:2611.09837}.

\bibitem{bajajlin}
K. Bajaj and H. W. Lin.
\newblock Solving the $\SU(2)$ BFSS model.
\newblock Preprint, November 2026.
\newblock \href{https://arxiv.org/abs/2611.14512}{arXiv:2611.14512}.

\bibitem{gwz}
H. Gluck, F. Warner, and W. Ziller.
\newblock The geometry of the Hopf fibrations.
\newblock \emph{L'Enseignement Mathématique} \textbf{32} (1986), 173--198.
\newblock \href{https://ncatlab.org/nlab/files/GluckWarnerZiller-HopfFibrations.pdf}{Article scan}.

\bibitem{pdg}
F. Takahashi et al. (Particle Data Group).
\newblock Review of Particle Physics.
\newblock \emph{International Journal of Modern Physics A} \textbf{41} (2026), 2630011.
\newblock \href{https://pdg.lbl.gov/2026/reviews/standard_model_and_related.html}{Standard Model and related reviews}.
\end{thebibliography}

\end{document}
